\section{OMP Coefficient Recovery and Base-Side Reconstruction}
\label{app:omp-coefficient-recovery}

This section supports the coefficient-reuse claim in Section~\ref{sec:methods}: if the attacker writes the donor row as a sparse mixture over donor anchors, then an OMP-based shared-basis transplant recovers the planted support and reuses nearly the planted coefficients on the base anchors. The proof uses the standard mutual-coherence exact-recovery condition for OMP, specialized to the donor composite dictionary \citep{tropp2004greed,troppgilbert2007omp,donohoeladtemlyakov2006stable,caiwang2011omp}.

\begin{assumption}[Normalized donor dictionary and mutual coherence]
\label{ass:omp-coherence}
Write $d_j := (\phi_j^{(d)})^\top \in \mathbb{R}^{\delta_d}$ for the donor anchor atom indexed by $j \in \mathcal{T}$, and let $D := \Phi_d^\top$. Assume $\|d_j\|_2=1$ for all $j\in\mathcal{T}$. Define
\[
\mu_d := \max_{i\neq j,\ i,j\in\mathcal{T}} |\langle d_i,d_j\rangle|.
\]
For a support $S\subseteq\mathcal{T}$ with $|S|=k$, assume
\[
\mu_d < \frac{1}{2k-1},
\qquad
\theta_k := 1-(2k-1)\mu_d > 0.
\]
\end{assumption}

\begin{assumption}[Noisy sparse attacker row and coefficient margin]
\label{ass:omp-noisy-row}
The planted donor row has the form
\[
x_d = \Phi_{d,S}^\top \alpha + \varepsilon = D_S\alpha+\varepsilon,
\qquad
\|\varepsilon\|_2 \le \eta_\varepsilon,
\]
where $D_S:=\Phi_{d,S}^\top$. The designed coefficients obey the beta-min condition
\[
\min_{j\in S}|\alpha_j| \ge c_\alpha,
\qquad
c_\alpha > \frac{2\eta_\varepsilon}{\theta_k}.
\]
\end{assumption}

\begin{lemma}[Coherence implies the donor-side ERC]
\label{lem:coherence-erc}
Under Assumption~\ref{ass:omp-coherence},
\[
\eta_S := \max_{\ell\notin S} \|D_S^\dagger d_\ell\|_1
\le
\frac{k\mu_d}{1-(k-1)\mu_d}
< 1.
\]
\end{lemma}

\begin{proof}
Let $G_S:=D_S^\top D_S$. Since the donor atoms are normalized and have mutual coherence $\mu_d$, the diagonal entries of $G_S$ equal $1$, and the absolute value of each off-diagonal entry is at most $\mu_d$. Hence $G_S=I+E$ with
\[
\|E\|_{1,1}\le (k-1)\mu_d < 1,
\]
where $\|M\|_{1,1}$ is the induced $\ell_1$ operator norm. The Neumann-series bound gives
\[
\|G_S^{-1}\|_{1,1}
\le
\frac{1}{1-(k-1)\mu_d}.
\]
For any $\ell\notin S$,
\[
D_S^\dagger d_\ell=(D_S^\top D_S)^{-1}D_S^\top d_\ell=G_S^{-1}D_S^\top d_\ell.
\]
Moreover, $\|D_S^\top d_\ell\|_1 \le k\mu_d$. Therefore
\[
\|D_S^\dagger d_\ell\|_1
\le
\|G_S^{-1}\|_{1,1}\|D_S^\top d_\ell\|_1
\le
\frac{k\mu_d}{1-(k-1)\mu_d}.
\]
The final quantity is smaller than $1$ because $\mu_d < (2k-1)^{-1}$.
\end{proof}

\begin{theorem}[OMP coefficient recovery under donor-side perturbation]
\label{thm:omp-coefficient-recovery}
Under Assumptions~\ref{ass:omp-coherence} and~\ref{ass:omp-noisy-row}, OMP applied to $x_d$ for exactly $k$ iterations selects the support $S$. The least-squares coefficients $\beta$ recovered on this support satisfy
\[
\|\beta-\alpha\|_2
\le
\frac{\|\varepsilon\|_2}{\sigma_{\min}(\Phi_{d,S})}
\le
\frac{\eta_\varepsilon}{\sigma_{\min}(\Phi_{d,S})}.
\]
Thus one may take
\[
C_{\mathrm{rec}}=1,
\qquad
\varepsilon_{\mathrm{rec}}
:=
\frac{\eta_\varepsilon}{\sigma_{\min}(\Phi_{d,S})}.
\]
\end{theorem}

\begin{proof}
Let $\Lambda\subseteq S$ be the set of atoms selected after some OMP iterations, and assume inductively that every selected atom belongs to $S$. Let $P_\Lambda^\perp$ be the orthogonal projector onto the complement of $\operatorname{span}(D_\Lambda)$. The residual is
\[
r=P_\Lambda^\perp x_d
=
P_\Lambda^\perp D_S\alpha+P_\Lambda^\perp\varepsilon
=
z+n,
\]
where
\[
z:=P_\Lambda^\perp D_{S\setminus\Lambda}\alpha_{S\setminus\Lambda},
\qquad
n:=P_\Lambda^\perp\varepsilon.
\]
Since $P_\Lambda^\perp$ is non-expansive and each atom is normalized,
\[
|\langle d_j,n\rangle|
=
|\langle P_\Lambda^\perp d_j,\varepsilon\rangle|
\le
\|\varepsilon\|_2
\le
\eta_\varepsilon
\]
for every $j\in\mathcal{T}$.

Let $R:=S\setminus\Lambda$. The matrix $D_R^\top P_\Lambda^\perp D_R$ is the Schur complement of $D_\Lambda^\top D_\Lambda$ inside $D_S^\top D_S$. Hence its smallest eigenvalue is at least $\lambda_{\min}(D_S^\top D_S)$. Gershgorin's theorem gives
\[
\lambda_{\min}(D_S^\top D_S)
\ge
1-(k-1)\mu_d.
\]
Therefore, with
\[
M_{\mathrm{true}}:=\max_{j\in R}|\langle d_j,z\rangle|=\|D_R^\top z\|_\infty,
\]
we have
\[
M_{\mathrm{true}}
\ge
\frac{\|D_R^\top z\|_2}{\sqrt{|R|}}
\ge
\frac{(1-(k-1)\mu_d)\|\alpha_R\|_2}{\sqrt{|R|}}
\ge
(1-(k-1)\mu_d)c_\alpha.
\]
For any $\ell\notin S$, since $z\in\operatorname{span}(D_S)$,
\[
\langle d_\ell,z\rangle
=
\langle D_SD_S^\dagger d_\ell,z\rangle
=
\langle D_S^\dagger d_\ell,D_S^\top z\rangle.
\]
The entries of $D_S^\top z$ indexed by $\Lambda$ vanish because $z\perp\operatorname{span}(D_\Lambda)$. Lemma~\ref{lem:coherence-erc} then implies
\[
|\langle d_\ell,z\rangle|
\le
\|D_S^\dagger d_\ell\|_1\|D_S^\top z\|_\infty
\le
\eta_S M_{\mathrm{true}},
\]
where
\[
\eta_S \le \frac{k\mu_d}{1-(k-1)\mu_d}.
\]
Thus
\[
\max_{j\in R}|\langle d_j,r\rangle|
\ge
M_{\mathrm{true}}-\eta_\varepsilon
\]
and
\[
\max_{\ell\notin S}|\langle d_\ell,r\rangle|
\le
\eta_S M_{\mathrm{true}}+\eta_\varepsilon.
\]
OMP selects a remaining true atom whenever $(1-\eta_S)M_{\mathrm{true}}>2\eta_\varepsilon$. Using the bounds above,
\[
(1-\eta_S)M_{\mathrm{true}}
\ge
\frac{1-(2k-1)\mu_d}{1-(k-1)\mu_d}(1-(k-1)\mu_d)c_\alpha
=
\theta_k c_\alpha
>
2\eta_\varepsilon.
\]
Hence each iteration selects an element of $S$. After $k$ iterations the selected support is exactly $S$.

Once the support is recovered, OMP returns the least-squares solution on $D_S$:
\[
\beta=D_S^\dagger x_d=D_S^\dagger(D_S\alpha+\varepsilon)=\alpha+D_S^\dagger\varepsilon.
\]
Therefore
\[
\|\beta-\alpha\|_2
\le
\|D_S^\dagger\|_{\mathrm{op}}\|\varepsilon\|_2
=
\frac{\|\varepsilon\|_2}{\sigma_{\min}(D_S)}
=
\frac{\|\varepsilon\|_2}{\sigma_{\min}(\Phi_{d,S})}.
\]
\end{proof}

\begin{corollary}[Base-side reconstruction perturbation]
\label{cor:omp-base-reconstruction-error}
Under the hypotheses of Theorem~\ref{thm:omp-coefficient-recovery},
\[
\|\widehat{x}_b-\mu_{\text{base}}\|_2
\le
\|\Phi_{b,S}^\top\alpha-\mu_{\text{base}}\|_2
+
\|\Phi_b\|_{\mathrm{op}}\|\beta-\alpha\|_2.
\]
In particular,
\[
\|\widehat{x}_b-\mu_{\text{base}}\|_2
\le
\|\Phi_{b,S}^\top\alpha-\mu_{\text{base}}\|_2
+
\|\Phi_b\|_{\mathrm{op}}\varepsilon_{\mathrm{rec}}.
\]
\end{corollary}

\begin{proof}
Since the recovered support is $S$,
\[
\widehat{x}_b=\Phi_{b,S}^\top\beta.
\]
Adding and subtracting $\Phi_{b,S}^\top\alpha$ gives
\[
\|\widehat{x}_b-\mu_{\text{base}}\|_2
\le
\|\Phi_{b,S}^\top\alpha-\mu_{\text{base}}\|_2
+
\|\Phi_{b,S}^\top(\beta-\alpha)\|_2.
\]
The second term is bounded by
\[
\|\Phi_{b,S}^\top(\beta-\alpha)\|_2
\le
\|\Phi_{b,S}\|_{\mathrm{op}}\|\beta-\alpha\|_2
\le
\|\Phi_b\|_{\mathrm{op}}\|\beta-\alpha\|_2.
\]
Substituting Theorem~\ref{thm:omp-coefficient-recovery} proves the result.
\end{proof}

\noindent\textbf{Intuition.}
OMP is stable here because the donor row is not arbitrary: it is a sparse mixture over a low-coherence anchor support, plus a perturbation smaller than the coefficient margin. Once the victim recovers the same support, the only base-side error is the attacker's design residual plus the coefficient-recovery perturbation passed through $\Phi_b^\top$.

\section{Asymmetric Realizability of Base Salience and Donor Inertness}
\label{app:asymmetric-realizability}

This section formalizes the geometric mechanism behind the asymmetric-realizability plots in Section~\ref{sec:experiments}: the same coefficient vector can be suppressed by the donor projection $U\Phi_{d,S}^\top$ while activating a base target through $\Phi_{b,S}^\top$. The GSVD construction below gives the feasible set, its approximation gap, and the coefficient cost in the paired anchor coordinates used by the breaker-token designer.

\begin{assumption}[Fixed support geometry]
\label{ass:fixed-support-geometry}
Fix a support $S\subseteq\mathcal{T}$ with $|S|=k$. Define
\[
A:=\Phi_{b,S}^\top\in\mathbb{R}^{\delta_b\times k},
\qquad
C:=U\Phi_{d,S}^\top\in\mathbb{R}^{m\times k}.
\]
Assume the stacked matrix $[A^\top\ C^\top]^\top$ has full column rank after removing any common null directions. Let
\[
A=U_A\operatorname{diag}(a_i)X^{-1},
\qquad
C=U_C\operatorname{diag}(c_i)X^{-1}
\]
be a thin generalized singular value decomposition on this quotient, where $X$ is invertible, the nonzero columns $u_i$ of $U_A$ and $v_i$ of $U_C$ are orthonormal in their respective spaces, $a_i,c_i\ge 0$, and $a_i^2+c_i^2=1$. The generalized singular values are $\sigma_i=a_i/c_i$ when $c_i>0$, and $\sigma_i=+\infty$ when $c_i=0<a_i$.
\end{assumption}

\begin{theorem}[GSVD construction of asymmetric realizability]
\label{thm:exact-asymmetric-realizability}
Under Assumption~\ref{ass:fixed-support-geometry}, define
\[
\mathcal{I}_0:=\{i:c_i=0<a_i\},
\qquad
\mathcal{I}_+:=\{i:c_i>0,\ a_i>0\}.
\]
Decompose
\[
\mu_{\text{base}}
=
\sum_{i\in\mathcal{I}_0}\mu_i u_i
+
\sum_{i\in\mathcal{I}_+}\mu_i u_i
+
\mu_\perp,
\qquad
\mu_i:=\langle \mu_{\text{base}},u_i\rangle,
\]
where $\mu_\perp$ is orthogonal to the base-side GSVD directions with $a_i>0$. Then
\[
A\operatorname{Null}(C)=\operatorname{span}\{u_i:i\in\mathcal{I}_0\}.
\]
Consequently:
\begin{enumerate}[leftmargin=*]
\item Exact donor-stealthy realization exists, meaning there is $\alpha$ with $A\alpha=\mu_{\text{base}}$ and $C\alpha=0$, if and only if
\[
\sum_{i\in\mathcal{I}_+}\mu_i u_i+\mu_\perp=0.
\]
\item The optimal donor-stealthy base reconstruction is the orthogonal projection of $\mu_{\text{base}}$ onto $A\operatorname{Null}(C)$, and its exact error is
\[
\operatorname{dist}\bigl(\mu_{\text{base}},A\operatorname{Null}(C)\bigr)^2
=
\sum_{i\in\mathcal{I}_+}\mu_i^2+\|\mu_\perp\|_2^2.
\]
\item An explicit coefficient vector attaining this error is
\[
\alpha_0=X\theta_0,
\qquad
(\theta_0)_i=
\begin{cases}
\mu_i/a_i, & i\in\mathcal{I}_0,\\
0, & i\notin\mathcal{I}_0.
\end{cases}
\]
It satisfies $C\alpha_0=0$ and
\[
A\alpha_0=\sum_{i\in\mathcal{I}_0}\mu_i u_i.
\]
If $a_{\min,0}:=\min_{i\in\mathcal{I}_0}a_i>0$, then
\[
\|X^{-1}\alpha_0\|_2
\le
\frac{\|P_0\mu_{\text{base}}\|_2}{a_{\min,0}},
\qquad
\|\alpha_0\|_2
\le
\frac{\|X\|_{\mathrm{op}}\|P_0\mu_{\text{base}}\|_2}{a_{\min,0}},
\]
where $P_0$ is the projector onto $A\operatorname{Null}(C)$.
\end{enumerate}
\end{theorem}

\begin{proof}
Write $\theta=X^{-1}\alpha$. Then
\[
A\alpha=\sum_i a_i\theta_i u_i,
\qquad
C\alpha=\sum_i c_i\theta_i v_i.
\]
The condition $C\alpha=0$ is equivalent to $c_i\theta_i=0$ for every $i$, since the $v_i$ are orthonormal. Thus any donor-stealthy coefficient vector has $\theta_i=0$ whenever $c_i>0$. Applying $A$ leaves only directions with $c_i=0$ and $a_i>0$, so
\[
A\operatorname{Null}(C)=\operatorname{span}\{u_i:i\in\mathcal{I}_0\}.
\]
The exact-realizability and distance claims follow by orthogonally decomposing $\mu_{\text{base}}$ into the subspace $A\operatorname{Null}(C)$, its base-visible complement inside $\operatorname{range}(A)$, and the component $\mu_\perp$ outside $\operatorname{range}(A)$. The displayed $\alpha_0$ has $c_i(\theta_0)_i=0$ for every $i$, hence $C\alpha_0=0$, and
\[
A\alpha_0=\sum_{i\in\mathcal{I}_0}a_i\frac{\mu_i}{a_i}u_i
=
\sum_{i\in\mathcal{I}_0}\mu_i u_i.
\]
The norm bounds follow from
\[
\|X^{-1}\alpha_0\|_2^2
=
\sum_{i\in\mathcal{I}_0}\frac{\mu_i^2}{a_i^2}
\le
\frac{\|P_0\mu_{\text{base}}\|_2^2}{a_{\min,0}^2}
\]
and $\|\alpha_0\|_2\le \|X\|_{\mathrm{op}}\|X^{-1}\alpha_0\|_2$.
\end{proof}

\begin{corollary}[Generic exact feasibility in a random-anchor model]
\label{cor:random-anchor-realizability}
Let $N_C\in\mathbb{R}^{k\times q}$ have orthonormal columns spanning $\operatorname{Null}(C)$, where $q=\dim\operatorname{Null}(C)$. If $q\ge \delta_b$ and, conditional on $C$, the matrix $AN_C$ has a distribution with a density on $\mathbb{R}^{\delta_b\times q}$, then
\[
\operatorname{rank}(AN_C)=\delta_b
\]
with probability one. Therefore $A\operatorname{Null}(C)=\mathbb{R}^{\delta_b}$ and every $\mu_{\text{base}}\in\mathbb{R}^{\delta_b}$ is exactly donor-stealth realizable. If, more specifically, $AN_C$ is an isotropic sub-Gaussian $\delta_b\times q$ matrix with scale $\sigma_A$, then for a universal constant $c>0$,
\[
\Pr\left[\sigma_{\min}(AN_C)\ge c\sigma_A(\sqrt{q}-\sqrt{\delta_b}-t)\right]
\ge
1-2\exp(-t^2)
\]
for every $0<t<\sqrt{q}-\sqrt{\delta_b}$ \citep{rudelsonvershynin2009smallest}.
\end{corollary}

\begin{proof}
The determinant of any $\delta_b\times\delta_b$ minor of $AN_C$ is a polynomial in the entries of $AN_C$. Because $q\ge \delta_b$, at least one such polynomial is not identically zero. A matrix distribution with a density assigns probability zero to the zero set of any nonzero polynomial, so $AN_C$ has full row rank with probability one. Then $AN_C\mathbb{R}^q=\mathbb{R}^{\delta_b}$, which is exactly $A\operatorname{Null}(C)=\mathbb{R}^{\delta_b}$. The sub-Gaussian bound is the standard smallest-singular-value estimate for rectangular random matrices applied to $AN_C$.
\end{proof}

\begin{lemma}[GSVD form of the regularized designer]
\label{thm:lambda-regularized-realizability}
Under the orthogonal coordinate model of Assumption~\ref{ass:orthogonal-gsvd} below, the solution of
\[
\alpha_\lambda
=
\arg\min_\alpha
\|A\alpha-\mu_{\text{base}}\|_2^2
+
\lambda\|C\alpha\|_2^2
+
\rho\|\alpha\|_2^2
\]
is
\[
\alpha_\lambda
=
\sum_i
\frac{a_i\mu_i}{a_i^2+\lambda c_i^2+\rho}z_i.
\]
The corresponding geometric error is
\[
\|A\alpha_\lambda-\mu_{\text{base}}\|_2^2
+
\lambda\|C\alpha_\lambda\|_2^2
=
\|\mu_\perp\|_2^2
+
\sum_i
\frac{(\lambda c_i^2+\rho)^2+\lambda a_i^2c_i^2}{(a_i^2+\lambda c_i^2+\rho)^2}\mu_i^2.
\]
\end{lemma}

\begin{proof}
Under Assumption~\ref{ass:orthogonal-gsvd}, write $\alpha=\sum_i t_i z_i$ and $\mu_{\text{base}}=\sum_i\mu_i u_i+\mu_\perp$. Orthogonality gives
\[
\|A\alpha-\mu_{\text{base}}\|_2^2
+
\lambda\|C\alpha\|_2^2
+
\rho\|\alpha\|_2^2
=
\|\mu_\perp\|_2^2+
\sum_i\left((a_it_i-\mu_i)^2+(\lambda c_i^2+\rho)t_i^2\right).
\]
Minimizing the one-dimensional quadratic in $t_i$ yields
\[
t_i=\frac{a_i\mu_i}{a_i^2+\lambda c_i^2+\rho}.
\]
Substituting this value gives the displayed error decomposition.
\end{proof}

\begin{remark}[Boundary cases]
\label{thm:asymmetric-failure-conditions}
The GSVD theorem identifies the only hard failures. If $\mu_{\text{base}}\perp A\operatorname{Null}(C)$, every exactly donor-stealthy coefficient vector satisfies $\|A\alpha-\mu_{\text{base}}\|_2\ge \|\mu_{\text{base}}\|_2$. If $\dim\operatorname{Null}(C)=0$, donor stealth forces $\alpha=0$, so the same lower bound holds with equality. These statements are boundary readings of Theorem~\ref{thm:exact-asymmetric-realizability}, not separate feasibility principles.
\end{remark}

\noindent\textbf{Intuition.}
The donor penalty removes GSVD coefficient directions with $c_i>0$. Exact asymmetric realization is possible precisely when the target lives in the base image of the remaining directions $c_i=0$. When exact stealth is impossible, the unavoidable base error is the energy of $\mu_{\text{base}}$ outside those donor-null directions, and the construction above gives the attacker coefficients that attain it.

\noindent\textbf{The base target $\mu_{\text{base}}$ is a presentation choice.}
Throughout the paper we instantiate the base-side target $\mu_{\text{base}}$ as the empirical mean of the base model's last-layer hidden states on a public WikiText sample (Section~\ref{subsec:pipeline}). This is a concrete, reproducible choice for the demonstrations and gives the breaker a generic, prompt-invariant high-salience direction. The realizability statement above does not depend on this specific target: Theorem~\ref{thm:exact-asymmetric-realizability} characterizes the entire feasibility cone of base-side directions admissible under donor stealth, and any direction in that cone supports the same construction. The centroid choice can be replaced with task-specific or token-specific base directions without altering the structural argument.

\section{A Wide $\lambda$-Window for Donor Suppression and Base Activation}
\label{app:lambda-window}

This section supports the $\lambda$-ablation claim documented in Appendix~\ref{app:lambda_window}: increasing the donor-suppression penalty can rapidly reduce donor activation while preserving base-side alignment. The statement is geometric; the conversion from geometry to SER is handled later in Theorem~\ref{thm:end-to-end-omp-attack}.

\begin{assumption}[Orthogonal generalized spectral coordinates]
\label{ass:orthogonal-gsvd}
For the fixed support $S$, let $A$ and $C$ be as in Assumption~\ref{ass:fixed-support-geometry}. Assume there exist orthonormal vectors $z_i\in\mathbb{R}^{k}$, orthonormal base-side vectors $u_i\in\mathbb{R}^{\delta_b}$, and orthonormal donor-side vectors $v_i\in\mathbb{R}^{m}$, with scalars $a_i,c_i\ge 0$, such that
\[
Az_i=a_i u_i,
\qquad
Cz_i=c_i v_i,
\qquad
a_i^2+c_i^2=1.
\]
For $c_i>0$, define $\sigma_i:=a_i/c_i$; if $c_i=0<a_i$, set $\sigma_i=+\infty$. Order the generalized singular values in non-increasing order, with the $+\infty$ directions placed first and the finite values ordered afterward.
\end{assumption}

\begin{theorem}[Spectral $\lambda$-window for simultaneous base alignment and donor suppression]
\label{thm:wide-lambda-window}
Assume Assumption~\ref{ass:orthogonal-gsvd}. Write
\[
\mu_{\text{base}}=\sum_i \mu_i u_i+\mu_\perp,
\qquad
\mu_i:=\langle \mu_{\text{base}},u_i\rangle,
\qquad
\mu_\perp\perp\operatorname{span}\{u_i\}_i.
\]
For $\lambda\ge 0$ and $\rho\ge 0$, let $\alpha_\lambda$ be the solution in Lemma~\ref{thm:lambda-regularized-realizability}. Then:
\begin{enumerate}[leftmargin=*]
\item The donor-visible norm obeys
\[
\|C\alpha_\lambda\|_2
\le
\frac{K_C}{\lambda}
\qquad
\text{for all }\lambda>0,
\]
where
\[
K_C:=\left(\sum_{i:c_i>0}\frac{a_i^2\mu_i^2}{c_i^2}\right)^{1/2}
=
\left(\sum_{i:c_i>0}\sigma_i^2\mu_i^2\right)^{1/2}.
\]
\item Fix a retained index set $\{1,\ldots,r\}$ among the ordered generalized singular directions and define
\[
E_r^2:=\|\mu_\perp\|_2^2+
\sum_{i>r}\mu_i^2,
\qquad
M_r^2:=\sum_{i=1}^{r}\mu_i^2.
\]
Assume $M_r>0$. Let $\varepsilon_{\text{base}}>E_r$ and set
\[
\Delta_r:=\left(\frac{\varepsilon_{\text{base}}^2-E_r^2}{M_r^2}\right)^{1/2},
\qquad
0<\Delta_r<1.
\]
For retained exact donor-stealth directions with $c_i=0$, assume
\[
\frac{\rho}{1+\rho}\le \Delta_r.
\]
For a donor target $\varepsilon_{\text{donor}}>0$, define
\[
\lambda_{\min}:=\frac{K_C}{\varepsilon_{\text{donor}}}
\]
and
\[
\lambda_{\max}:=
\min_{\substack{1\le i\le r\\ c_i>0}}
\left\{
\frac{\Delta_r}{1-\Delta_r}\sigma_i^2-\frac{\rho}{c_i^2}
\right\},
\]
with $\lambda_{\max}=+\infty$ if all retained directions have $c_i=0$. If $\lambda_{\min}\le \lambda\le \lambda_{\max}$, then
\[
\|C\alpha_\lambda\|_2\le \varepsilon_{\text{donor}},
\qquad
\|A\alpha_\lambda-\mu_{\text{base}}\|_2\le \varepsilon_{\text{base}}.
\]
\item Let
\[
g_r:=\frac{\sigma_r}{\sigma_{r+1}}
\]
for finite $\sigma_r,\sigma_{r+1}$, and let $E_{r,\sigma}^2:=\sum_{i>r}\mu_i^2$. Since
\[
K_C
\le
\sigma_r\left(\frac{\sigma_1}{\sigma_r}M_r+\frac{1}{g_r}E_{r,\sigma}\right),
\]
whenever the interval is nonempty its certified width satisfies
\[
\lambda_{\max}-\lambda_{\min}
\ge
\frac{\Delta_r}{1-\Delta_r}\sigma_r^2
-
\max_{\substack{1\le i\le r\\ c_i>0}}\frac{\rho}{c_i^2}
-
\frac{\sigma_r}{\varepsilon_{\text{donor}}}
\left(\frac{\sigma_1}{\sigma_r}M_r+\frac{1}{g_r}E_{r,\sigma}\right).
\]
Thus the prediction is falsifiable from the anchor pair: retained base energy $M_r$, tolerated tail $E_r$, and the spectral gap $g_r$ determine whether a broad interval of $\lambda$ values should preserve base alignment while suppressing donor projection.
\end{enumerate}
\end{theorem}

\begin{proof}
The formula for $\alpha_\lambda$ follows from Lemma~\ref{thm:lambda-regularized-realizability}. For the donor norm,
\[
C\alpha_\lambda
=
\sum_i
\frac{a_ic_i\mu_i}{a_i^2+\lambda c_i^2+\rho}v_i.
\]
When $c_i=0$, the term vanishes. When $c_i>0$,
\[
\left|
\frac{a_ic_i\mu_i}{a_i^2+\lambda c_i^2+\rho}
\right|
\le
\frac{|a_i\mu_i|}{\lambda c_i}.
\]
Squaring and summing gives $\|C\alpha_\lambda\|_2\le K_C/\lambda$.

The base residual is
\[
A\alpha_\lambda-\mu_{\text{base}}
=
-
\sum_i
\frac{\lambda c_i^2+\rho}{a_i^2+\lambda c_i^2+\rho}\mu_i u_i
-
\mu_\perp.
\]
For $i\le r$ with $c_i>0$, the shrinkage factor is at most $\Delta_r$ precisely when
\[
\lambda
\le
\frac{\Delta_r}{1-\Delta_r}\frac{a_i^2}{c_i^2}-\frac{\rho}{c_i^2}
=
\frac{\Delta_r}{1-\Delta_r}\sigma_i^2-\frac{\rho}{c_i^2}.
\]
For $i\le r$ with $c_i=0$, Assumption~\ref{ass:orthogonal-gsvd} gives $a_i=1$, and the shrinkage factor equals $\rho/(1+\rho)$, which is at most $\Delta_r$ by the theorem hypothesis. For $i>r$, the shrinkage factor is at most $1$. Therefore $\lambda\le\lambda_{\max}$ implies
\[
\|A\alpha_\lambda-\mu_{\text{base}}\|_2^2
\le
\Delta_r^2M_r^2+E_r^2
=
\varepsilon_{\text{base}}^2.
\]
If $\lambda\ge\lambda_{\min}=K_C/\varepsilon_{\text{donor}}$, the donor norm is at most $\varepsilon_{\text{donor}}$. This proves the window.

For the width bound, use $\sigma_i\ge\sigma_r$ for $i\le r$ to lower bound $\lambda_{\max}$ and split $K_C$ into retained and discarded parts:
\[
K_C^2
=
\sum_{i\le r}\sigma_i^2\mu_i^2+
\sum_{i>r}\sigma_i^2\mu_i^2
\le
\sigma_1^2M_r^2+\sigma_{r+1}^2E_{r,\sigma}^2.
\]
The displayed upper bound on $K_C$ follows from $\sqrt{x^2+y^2}\le x+y$ and $\sigma_{r+1}=\sigma_r/g_r$. Subtracting $\lambda_{\min}$ from the lower bound on $\lambda_{\max}$ completes the proof.
\end{proof}

\noindent\textbf{Intuition.}
Large generalized singular values are directions that the base sees much more strongly than the donor. Increasing $\lambda$ suppresses donor-visible components at rate $1/\lambda$, while high-$\sigma_i$ base components remain accurate until $\lambda$ reaches their squared spectral scale. The experimentally observed window $\lambda\in[768,1536]$ is the regime where the measured pair has a large retained spectral scale and a small discarded tail.

\section{From Base Reconstruction to Sequence Emission Rate}
\label{app:base-reconstruction-to-ser}

This section connects geometric reconstruction error to the SER metric of Section~\ref{sec:methods}. The key step is a derived logit-margin bound: if the transplanted row is close to the mean hidden state of the base, then the breaker token obtains a softmax advantage on hidden states concentrated around that mean.

\begin{assumption}[Measurable hidden-state and competitor-row controls]
\label{ass:hidden-norm-control}
Let $h_t$ be the last-layer hidden state of the post-transplant base at a generation position $t$, under the prompt distribution and under histories that have not yet emitted $\tau$. Define
\[
\mu_h:=\mathbb{E}[h_t].
\]
Assume that for every evaluated step and every no-prior-$\tau$ history,
\[
\Pr\left[\|h_t-\mu_h\|_2\le \xi_b\mid \text{history}\right]
\ge
1-\zeta_b.
\]
Let $R_b:=\max_{v\neq\tau}\|w_v\|_2$, and let $V_b$ be the base vocabulary size after transplant. These are directly measurable from the post-transplant base model.
\end{assumption}

\begin{lemma}[Logit margin from base-row alignment]
\label{lem:logit-margin-from-alignment}
Under Assumption~\ref{ass:hidden-norm-control}, suppose the breaker LM-head row satisfies
\[
\|w_\tau-\mu_h\|_2\le r_b.
\]
On the event $\|h_t-\mu_h\|_2\le \xi_b$,
\[
w_\tau^\top h_t-\max_{v\neq\tau}w_v^\top h_t
\ge
m_b(r_b),
\]
where
\[
m_b(r_b)
:=
\|\mu_h\|_2^2
-r_b\|\mu_h\|_2
-(\|\mu_h\|_2+r_b)\xi_b
-R_b(\|\mu_h\|_2+\xi_b).
\]
Consequently, if $m_b(r_b)>0$, the conditional softmax probability of $\tau$ on this event is at least
\[
q_b(r_b)
:=
\frac{1}{1+(V_b-1)\exp[-m_b(r_b)]}.
\]
\end{lemma}

\begin{proof}
Write $w_\tau=\mu_h+e$ with $\|e\|_2\le r_b$ and $h_t=\mu_h+\xi$ with $\|\xi\|_2\le \xi_b$. Then
\[
w_\tau^\top h_t
=
\|\mu_h\|_2^2+e^\top\mu_h+\mu_h^\top\xi+e^\top\xi
\ge
\|\mu_h\|_2^2-r_b\|\mu_h\|_2-(\|\mu_h\|_2+r_b)\xi_b.
\]
For every competitor $v\neq\tau$,
\[
w_v^\top h_t
\le
\|w_v\|_2\|h_t\|_2
\le
R_b(\|\mu_h\|_2+\xi_b).
\]
Subtracting yields the margin. If the breaker logit exceeds every competitor by $m$, then
\[
\operatorname{softmax}(Wh_t)_\tau
=
\frac{1}{1+\sum_{v\neq\tau}\exp(w_v^\top h_t-w_\tau^\top h_t)}
\ge
\frac{1}{1+(V_b-1)e^{-m}}.
\]
Setting $m=m_b(r_b)$ proves the result.
\end{proof}

\begin{remark}[SER compounding step]
\label{rem:ser-compounding-step}
If, at every step and conditional on any history that has not yet emitted $\tau$, the probability of sampling $\tau$ is at least $q$, then the probability that $\tau$ appears at least once in a $T$-token continuation is at least $1-(1-q)^T$. In Theorem~\ref{thm:end-to-end-omp-attack}, Lemma~\ref{lem:logit-margin-from-alignment} supplies $q=(1-\zeta_b)q_b(r_b)$ for the attacked base. The experiments use $T=256$, temperature $1.0$, and top-$p=0.9$; the same bound applies to top-$p$ sampling when $\tau$ remains in the retained nucleus on the margin event.
\end{remark}

\section{Operator-Specific Reachability for Differentiable Transplants}
\label{app:operator-specific-reachability}

This section refines the nonlinear-solver discussion in Appendix~\ref{app:nonlinear-shared-anchor}. OMP is handled by Theorem~\ref{thm:omp-coefficient-recovery}; here we prove that FOCUS softmax, top-$k$ mixtures, CLP, and WECHSEL each contain target-weight neighborhoods whenever their donor score maps can fit the corresponding local logits or positive scores.

\begin{theorem}[Shared-basis operator-class reachability]
\label{thm:operator-class-reachability}
Let $g_d:\mathbb{R}^{\delta_d}\to\mathcal{W}\subseteq\mathbb{R}^{N}$ be a donor-side coefficient map and let the base reconstruction be $x_b(x_d)=\Phi_b^\top g_d(x_d)$. Fix $w^\star\in\mathcal{W}$ and $\eta\ge 0$. If there exists a donor row $x_d$ such that
\[
\|g_d(x_d)-w^\star\|_1\le \eta,
\]
then
\[
\|x_b(x_d)-\Phi_b^\top w^\star\|_2
\le
\max_{j\in\mathcal{T}}\|\phi_j^{(b)}\|_2\eta.
\]
If instead the coefficient error is measured in $\ell_2$, then
\[
\|x_b(x_d)-\Phi_b^\top w^\star\|_2
\le
\|\Phi_b\|_{\mathrm{op}}\|g_d(x_d)-w^\star\|_2.
\]
\end{theorem}

\begin{proof}
The $\ell_1$ bound follows from
\[
\|\Phi_b^\top(g_d(x_d)-w^\star)\|_2
\le
\sum_{j\in\mathcal{T}}|g_{d,j}(x_d)-w_j^\star|\,\|\phi_j^{(b)}\|_2
\le
\max_{j\in\mathcal{T}}\|\phi_j^{(b)}\|_2\eta.
\]
The $\ell_2$ bound is the operator-norm inequality.
\end{proof}

\begin{assumption}[Softmax attention operator]
\label{ass:softmax-attention-operator}
The donor anchors and donor row are unit-normalized:
\[
\widehat{\phi}^{(d)}_j=\frac{\phi^{(d)}_j}{\|\phi^{(d)}_j\|_2},
\qquad
\widehat{x}_d=\frac{x_d}{\|x_d\|_2}.
\]
For an inverse temperature $\beta_{\mathrm{sm}}>0$, the FOCUS coefficient map is
\[
w_j(x_d)=\operatorname{softmax}_j\left(\beta_{\mathrm{sm}}\langle \widehat{x}_d,\widehat{\phi}^{(d)}_j\rangle\right).
\]
Equivalently, with normalized rows,
\[
w(x_d)=\operatorname{softmax}(\beta_{\mathrm{sm}}\Phi_d\widehat{x}_d).
\]
\end{assumption}

\begin{lemma}[FOCUS softmax reachability]
\label{lem:softmax-reachability}
Let $w^\star\in\mathbb{R}^{N}$ lie in the open simplex, and set $w^\star_{\min}:=\min_j w^\star_j>0$. For any $c\in\mathbb{R}$ define
\[
\ell^\star(c):=\frac{\log w^\star+c\mathbf{1}}{\beta_{\mathrm{sm}}}.
\]
If there exists a unit vector $z\in\mathbb{R}^{\delta_d}$ such that
\[
\|\Phi_d z-\ell^\star(c)\|_2\le \eta_{\mathrm{fit}},
\]
then any nonzero $x_d$ with $\widehat{x}_d=z$ satisfies
\[
\|\operatorname{softmax}(\beta_{\mathrm{sm}}\Phi_d\widehat{x}_d)-w^\star\|_1
\le
\frac{\beta_{\mathrm{sm}}\sqrt{N}}{2}\eta_{\mathrm{fit}}.
\]
The role of $w^\star_{\min}$ is to ensure finite target logits; the Lipschitz constant itself is uniform over the simplex.
\end{lemma}

\begin{proof}
Softmax is invariant to adding $c\mathbf{1}$ to the logits, so $\operatorname{softmax}(\log w^\star+c\mathbf{1})=w^\star$. The Jacobian of softmax at probability vector $\pi$ is $\operatorname{diag}(\pi)-\pi\pi^\top$, the covariance matrix of a categorical variable, and its $\ell_2$ operator norm is at most $1/2$. Therefore
\[
\|\operatorname{softmax}(s)-\operatorname{softmax}(s')\|_2
\le
\frac{1}{2}\|s-s'\|_2.
\]
Using $\|u\|_1\le\sqrt{N}\|u\|_2$ with $s=\beta_{\mathrm{sm}}\Phi_d z$ and $s'=\log w^\star+c\mathbf{1}$ gives the bound.
\end{proof}

\begin{assumption}[$k$-NN / top-$k$ softmax operator]
\label{ass:knn-topk-operator}
The operator first selects
\[
\mathcal{N}_k(x_d)=\operatorname{TopK}_{j\in\mathcal{T}}\langle \phi_j^{(d)},x_d\rangle
\]
and then applies a softmax rule restricted to $\mathcal{N}_k(x_d)$. The target support $S\subseteq\mathcal{T}$ has $|S|=k$.
\end{assumption}

\begin{lemma}[Margin-stable support selection]
\label{lem:margin-stable-knn}
Under Assumption~\ref{ass:knn-topk-operator}, suppose $x_d$ satisfies
\[
\min_{j\in S}\langle \phi_j^{(d)},x_d\rangle
>
\max_{\ell\notin S}\langle \phi_\ell^{(d)},x_d\rangle+\gamma
\]
for some $\gamma>0$. Then $\mathcal{N}_k(x_d)=S$. Moreover, if $x_d'$ satisfies
\[
\max_{j\in\mathcal{T}}
|\langle \phi_j^{(d)},x_d'-x_d\rangle|
\le
\frac{\gamma}{2},
\]
then $\mathcal{N}_k(x_d')=S$ as well.
\end{lemma}

\begin{proof}
The first claim follows because every score indexed by $S$ exceeds every score outside $S$, and $|S|=k$. For the perturbation claim, for $j\in S$ and $\ell\notin S$,
\[
\langle \phi_j^{(d)},x_d'\rangle
\ge
\langle \phi_j^{(d)},x_d\rangle-\frac{\gamma}{2}
>
\langle \phi_\ell^{(d)},x_d\rangle+\frac{\gamma}{2}
\ge
\langle \phi_\ell^{(d)},x_d'\rangle.
\]
Thus the top-$k$ set remains $S$.
\end{proof}

\begin{corollary}[Top-$k$ mixture reachability]
\label{cor:topk-mixture-reachability}
Let $w^\star$ be supported on $S$ and lie in the open simplex over $S$. If $x_d$ satisfies the margin condition in Lemma~\ref{lem:margin-stable-knn} and the restricted softmax logits over $S$ satisfy
\[
\left\|
\Phi_{d,S}\widehat{x}_d
-
\frac{\log w^\star_S+c\mathbf{1}_S}{\beta_{\mathrm{sm}}}
\right\|_2
\le
\eta_{\mathrm{fit},S},
\]
then the full top-$k$ mixture weights satisfy
\[
\|w(x_d)-w^\star\|_1
\le
\frac{\beta_{\mathrm{sm}}\sqrt{k}}{2}\eta_{\mathrm{fit},S}.
\]
\end{corollary}

\begin{proof}
By Lemma~\ref{lem:margin-stable-knn}, the selected support is $S$. Outside $S$, both the target and recovered weights are zero. Inside $S$, apply Lemma~\ref{lem:softmax-reachability} with $N$ replaced by $k$.
\end{proof}

\begin{assumption}[CLP ReLU-renormalized coefficient map]
\label{ass:clp-operator}
For normalized donor row $\widehat{x}_d$, define scores
\[
s_j(x_d):=\langle \widehat{x}_d,\widehat{\phi}^{(d)}_j\rangle.
\]
The CLP coefficient map is
\[
w_j(x_d)=
\frac{\operatorname{ReLU}(s_j(x_d))}{\epsilon_{\mathrm{clp}}+\sum_{\ell\in\mathcal{T}}\operatorname{ReLU}(s_\ell(x_d))},
\]
with $\epsilon_{\mathrm{clp}}\ge 0$.
\end{assumption}

\begin{lemma}[CLP reachability through positive score fitting]
\label{lem:clp-reachability}
Let $w^\star$ lie in the simplex over a support $S$. Suppose there are $a>0$ and $x_d$ such that, with
\[
r_j:=\operatorname{ReLU}(s_j(x_d)),
\]
the positive score vector satisfies
\[
\|r-a w^\star\|_1\le \eta_{\mathrm{clp}},
\qquad
0\le \eta_{\mathrm{clp}}<a.
\]
Then the CLP weights satisfy
\[
\|w(x_d)-w^\star\|_1
\le
\frac{2\eta_{\mathrm{clp}}}{a-\eta_{\mathrm{clp}}}
+
\frac{\epsilon_{\mathrm{clp}}}{\epsilon_{\mathrm{clp}}+a-\eta_{\mathrm{clp}}}.
\]
\end{lemma}

\begin{proof}
Let $R:=\sum_j r_j$. Since $\sum_j w_j^\star=1$,
\[
|R-a|\le \|r-aw^\star\|_1\le \eta_{\mathrm{clp}},
\]
so $R\ge a-\eta_{\mathrm{clp}}>0$. Decompose
\[
\left\|\frac{r}{\epsilon_{\mathrm{clp}}+R}-w^\star\right\|_1
\le
\left\|\frac{r}{\epsilon_{\mathrm{clp}}+R}-\frac{r}{R}\right\|_1
+
\left\|\frac{r}{R}-w^\star\right\|_1.
\]
The first term equals
\[
\frac{\epsilon_{\mathrm{clp}}}{\epsilon_{\mathrm{clp}}+R}
\le
\frac{\epsilon_{\mathrm{clp}}}{\epsilon_{\mathrm{clp}}+a-\eta_{\mathrm{clp}}}.
\]
For the second term,
\[
\frac{r}{R}-w^\star
=
\frac{r-aw^\star}{R}+
\left(\frac{a}{R}-1\right)w^\star,
\]
so
\[
\left\|\frac{r}{R}-w^\star\right\|_1
\le
\frac{\eta_{\mathrm{clp}}}{R}+\frac{|a-R|}{R}
\le
\frac{2\eta_{\mathrm{clp}}}{a-\eta_{\mathrm{clp}}}.
\]
\end{proof}

\begin{assumption}[WECHSEL affine alignment]
\label{ass:wechsel-affine}
The WECHSEL operator uses an affine alignment
\[
\widetilde{x}(x_d)=A_{\mathrm{WEC}}(x_d-\mu_{\mathrm{src}})+\mu_{\mathrm{tgt}},
\]
followed by local $k$-NN softmax interpolation:
\[
f_{\mathrm{WEC}}(x_d)=
\sum_{v\in\mathcal{N}_k(\widetilde{x}(x_d))}
w_v(\widetilde{x}(x_d))\phi_v^{(b)}.
\]
Assume $A_{\mathrm{WEC}}$ has a right inverse on the target affine subspace: for every target $z$ in that subspace,
\[
x_d(z)=\mu_{\mathrm{src}}+A_{\mathrm{WEC}}^\dagger(z-\mu_{\mathrm{tgt}})
\]
satisfies $\widetilde{x}(x_d(z))=z$.
\end{assumption}

\begin{lemma}[WECHSEL reachability decomposition]
\label{lem:wechsel-reachability}
Under Assumption~\ref{ass:wechsel-affine},
\[
\|\widetilde{x}(x_d)-\widetilde{x}(x_d')\|_2
\le
\|A_{\mathrm{WEC}}\|_{\mathrm{op}}\|x_d-x_d'\|_2,
\]
and, for reachable targets $z,z'$,
\[
\|x_d(z)-x_d(z')\|_2
\le
\|A_{\mathrm{WEC}}^\dagger\|_{\mathrm{op}}\|z-z'\|_2.
\]
Consequently, any target aligned point $z$ satisfying the support margin condition of Lemma~\ref{lem:margin-stable-knn} and the restricted softmax fitting condition of Corollary~\ref{cor:topk-mixture-reachability} is realized by $x_d(z)$. The reachable WECHSEL weight family is the inclusion
\[
\mathcal{W}_{\mathrm{WEC}}
\supseteq
\bigcup_{\substack{S\subseteq\mathcal{T}\\ |S|=k}}
\left\{
 w\in\Delta_S:
 \exists z\in\operatorname{Im}(\widetilde{x})\text{ satisfying the }S\text{-margin condition and the }S\text{-softmax fit for }w
\right\}.
\]
If WECHSEL selects supports only by this margin rule and weights only by the stated local softmax, the inclusion is equality up to score ties.
\end{lemma}

\begin{proof}
For any $x_d,x_d'$,
\[
\widetilde{x}(x_d)-\widetilde{x}(x_d')=A_{\mathrm{WEC}}(x_d-x_d'),
\]
which gives the first Lipschitz bound. For reachable targets,
\[
x_d(z)-x_d(z')=A_{\mathrm{WEC}}^\dagger(z-z'),
\]
which gives the second bound. If $z$ satisfies the margin and fitting conditions, $x_d(z)$ maps to $z$, the selected support is $S$ by Lemma~\ref{lem:margin-stable-knn}, and the weights on $S$ are controlled by Corollary~\ref{cor:topk-mixture-reachability}. The displayed family follows.
\end{proof}

\noindent\textbf{Intuition.}
All differentiable operators have the same core: fit an operator-specific donor score vector, obtain weights, and reuse those weights on base anchors. FOCUS fits log-probabilities, top-$k$ methods add a margin-stable support, CLP fits positive ReLU scores, and WECHSEL first moves the donor row through an affine alignment before the same local-mixture argument applies.

\section{A Constructive Lower Bound for Single-Statistic Spectral Detectors}
\label{app:spectral-detector-lower-bound}

This section formalizes the spectral-mimicry claim in Section~\ref{sec:mimicry}. The result applies to any detector that thresholds one scalar donor-row statistic, once the feasible design set contains a controllable direction for that statistic. The proof has two parts: the base-alignment and donor-suppression constraints leave an affine design set in which the attacker can match the natural pushforward of common scalar statistics; once the pushforward distributions are close, a single-statistic threshold cannot catch attacker rows without flagging natural tokens.

\begin{assumption}[Single-statistic threshold detector]
\label{ass:single-statistic-threshold}
Let $s:\mathbb{R}^{\delta_d}\to\mathbb{R}$ be a fixed scalar statistic of a donor embedding row. A detector is a measurable threshold rule
\[
D(x)=\mathbf{1}\{s(x)\in B\},
\]
where $B\subseteq\mathbb{R}$ is the rejection region. Let $\nu_{\mathrm{nat}}$ be the empirical distribution of $s(x)$ over natural donor tokens.
\end{assumption}

\begin{assumption}[Design freedom after base and donor constraints]
\label{ass:detector-design-freedom}
Let
\[
\mathcal{F}_{\alpha}
:=
\{\alpha:A\alpha=\mu_{\text{base}},\ C\alpha=0\}
\]
be nonempty, and fix one solution $\alpha_0$. Then
\[
\mathcal{F}_{\alpha}=\alpha_0+\operatorname{Null}\left(\begin{bmatrix}A\\ C\end{bmatrix}\right).
\]
Equivalently, the donor rows available without changing the exact base image or donor projection contain the affine set
\[
\mathcal{F}_{x}
:=
x_0+L,
\qquad
x_0:=\Phi_{d,S}^\top\alpha_0,
\qquad
L:=\Phi_{d,S}^\top\operatorname{Null}\left(\begin{bmatrix}A\\ C\end{bmatrix}\right).
\]
Let $d_f:=\dim L$. For scalar base constraints such as a fixed alignment $\langle A\alpha,\mu_{\text{base}}\rangle$, the same calculation gives $d_f\ge \dim\operatorname{Null}(C)-1$ before accounting for degeneracy of $\Phi_{d,S}^\top$; for exact vector constraints, $d_f\ge \dim\operatorname{Null}(C)-\operatorname{rank}(A|_{\operatorname{Null}(C)})$ before the same donor-map quotient.
\end{assumption}

\begin{assumption}[Statistic controllability on the design set]
\label{ass:statistic-controllability}
There is an interval $I\subseteq\mathbb{R}$ and a measurable map $\gamma:I\to\mathcal{F}_x$ such that
\[
s(\gamma(t))=t
\qquad
\text{for every }t\in I.
\]
Let $\epsilon_{\mathrm{tail}}:=\nu_{\mathrm{nat}}(I^c)$.
\end{assumption}

\begin{theorem}[Constructive mimicry and threshold-detector lower bound]
\label{thm:spectral-threshold-lower-bound}
Under Assumptions~\ref{ass:single-statistic-threshold}--\ref{ass:statistic-controllability}, there is a distribution $\nu_{\mathrm{atk}}$ over attacker-designed donor rows in $\mathcal{F}_x$ such that
\[
\operatorname{TV}\left(s_\#\nu_{\mathrm{atk}},\nu_{\mathrm{nat}}\right)
\le
\epsilon_{\mathrm{tail}}.
\]
Consequently, for any threshold detector $D(x)=\mathbf{1}\{s(x)\in B\}$, if
\[
\operatorname{FNR}_{\mathrm{atk}}:=s_\#\nu_{\mathrm{atk}}(B^c)\le \delta,
\]
then
\[
\operatorname{FPR}_{\mathrm{nat}}:=\nu_{\mathrm{nat}}(B)
\ge
1-\delta-\epsilon_{\mathrm{tail}}.
\]
In particular, when $\delta\le 1/2$, $\operatorname{FPR}_{\mathrm{nat}}\ge \delta-\epsilon_{\mathrm{tail}}$.
\end{theorem}

\begin{proof}
Let $\bar\nu$ be the conditional distribution of the natural statistic given $I$:
\[
\bar\nu(E)=\frac{\nu_{\mathrm{nat}}(E\cap I)}{\nu_{\mathrm{nat}}(I)}.
\]
Draw $T\sim\bar\nu$ and output the attacker row $X=\gamma(T)$. By Assumption~\ref{ass:statistic-controllability}, $s(X)=T$, so $s_\#\nu_{\mathrm{atk}}=\bar\nu$. The total variation distance between a distribution and its conditioning on an event of mass $1-\epsilon_{\mathrm{tail}}$ is at most $\epsilon_{\mathrm{tail}}$, proving the constructive TV bound.

For the detector lower bound, $\operatorname{FNR}_{\mathrm{atk}}\le\delta$ implies $s_\#\nu_{\mathrm{atk}}(B)\ge 1-\delta$. Total variation gives
\[
\nu_{\mathrm{nat}}(B)
\ge
s_\#\nu_{\mathrm{atk}}(B)-\epsilon_{\mathrm{tail}}
\ge
1-\delta-\epsilon_{\mathrm{tail}}.
\]
The final weaker bound follows from $1-\delta\ge\delta$ when $\delta\le 1/2$.
\end{proof}

\begin{corollary}[Controllability for common donor-row statistics]
\label{cor:common-statistic-controllability}
Assume $d_f\ge 1$ and the feasible row set contains $x_0+tq$ for $|t|\le R$, where $q\in L$ is a unit vector.
\begin{enumerate}[leftmargin=*]
\item For the row-norm statistic $s(x)=\|x\|_2$, if $q\perp x_0$, then $s(x_0+tq)$ covers $[\|x_0\|_2,\sqrt{\|x_0\|_2^2+R^2}]$.
\item For the donor PCA residual statistic $s_U(x)=\|(I-U^\top U)x\|_2$, if $q\in\operatorname{Null}(U)$ and $q\perp (I-U^\top U)x_0$, then $s_U(x_0+tq)$ covers $[\|(I-U^\top U)x_0\|_2,\sqrt{\|(I-U^\top U)x_0\|_2^2+R^2}]$.
\item More generally, any Lipschitz-1 scalar statistic $s$ satisfying $s(x_0+t_-q)\le a$ and $s(x_0+t_+q)\ge b$ for some $t_-,t_+\in[-R,R]$ is controllable on $I=[a,b]$ by the intermediate value theorem after restricting to a monotone measurable branch.
\end{enumerate}
Thus norm detectors, PCA-residual detectors of the Magikarp type \citep{land2024fishingmagikarpautomaticallydetecting}, and any scalar Lipschitz statistic with a controllable direction fall under Theorem~\ref{thm:spectral-threshold-lower-bound}.
\end{corollary}

\begin{proof}
The first two claims follow by direct calculation:
\[
\|x_0+tq\|_2^2=\|x_0\|_2^2+t^2
\]
when $q\perp x_0$, and similarly
\[
\|(I-U^\top U)(x_0+tq)\|_2^2
=
\|(I-U^\top U)x_0\|_2^2+t^2
\]
when $q\in\operatorname{Null}(U)$ and $q\perp (I-U^\top U)x_0$. The third claim uses continuity of Lipschitz functions on the line segment and the intermediate value theorem.
\end{proof}

\begin{remark}[Finite multi-statistic extension]
\label{rem:multi-statistic-detectors}
For a vector statistic $s=(s_1,\ldots,s_p)$, the same proof applies with $I\subseteq\mathbb{R}^p$ whenever the feasible design set admits a measurable right inverse $\gamma:I\to\mathcal{F}_x$ with $s(\gamma(t))=t$. A sufficient local condition is that $s$ is continuously differentiable and its Jacobian restricted to $L$ has rank $p$ at some feasible row; the inverse function theorem then gives controllability on a neighborhood. The TV lower bound is unchanged except for the natural mass outside that controllable neighborhood.
\end{remark}

\noindent\textbf{Intuition.}
After base alignment and donor suppression are fixed, the attacker often retains an affine set of coefficient perturbations that do not change either constraint. If a detector only reads one scalar statistic, one controllable direction in this affine set is enough to match the natural statistic distribution on its central support; thresholding that statistic then trades false negatives directly against false positives.

\section{Calibration-Based Recipient-Side SER Bound}
\label{app:recipient-side-defense-ser}

This section turns the differential-SER auditing protocol in Appendix~\ref{app:auditing} into a sufficient condition for bounded sampling SER. It applies to recipient-side checks that reject newly synthesized rows whose calibration probability is too often above a chosen per-step threshold.

\begin{assumption}[Calibration transfer for synthesized rows]
\label{ass:calibration-transfer}
For a newly synthesized embedding or LM-head row $w_\tau$, define
\[
\widehat{q}_\tau
:=
\Pr_{p\sim\mathcal{P}_{\mathrm{cal}}}
\left[
\operatorname{softmax}(Wh_p)_\tau\ge q^\star
\right].
\]
Assume the deployment prompt-position distribution is dominated by the calibration distribution in the following measurable sense: if $\widehat{q}_\tau<\delta$, then for every generation position $t\in\{1,\ldots,T\}$,
\[
\Pr\left[
\operatorname{softmax}(Wh_t)_\tau\ge q^\star
\right]
\le
\delta,
\]
where the probability is over deployment prompts and decoding histories before position $t$.
\end{assumption}

\begin{theorem}[Calibration-based SER bound]
\label{thm:calibration-ser-bound}
Suppose the recipient-side defense rejects $w_\tau$ whenever $\widehat{q}_\tau\ge\delta$. Under Assumption~\ref{ass:calibration-transfer}, any token $\tau$ that passes the check satisfies
\[
\operatorname{SER}(\tau)
\le
1-(1-q^\star)^T+
\delta T
\]
under temperature-$1$ sampling, including top-$p$ sampling when the post-truncation probability of $\tau$ is no larger than its checked probability on the non-tail event.
\end{theorem}

\begin{proof}
Let
\[
B_t:=\left\{\operatorname{softmax}(Wh_t)_\tau\ge q^\star\right\}.
\]
Since the token passes the check, $\widehat{q}_\tau<\delta$, and Assumption~\ref{ass:calibration-transfer} gives $\Pr(B_t)\le\delta$ for every $t$. By the union bound,
\[
\Pr\left(\bigcup_{t=1}^{T}B_t\right)
\le
\delta T.
\]
On the complement of $\bigcup_t B_t$, the probability of sampling $\tau$ at every position is at most $q^\star$. Iterating the conditional probabilities along the generation history gives probability at least $(1-q^\star)^T$ of never sampling $\tau$ on this non-tail event. Hence, conditional on no tail event, the probability of at least one occurrence is at most $1-(1-q^\star)^T$. Adding the tail probability gives the bound.
\end{proof}

\noindent\textbf{Intuition.}
The audit bounds the mass of high-probability positions for a synthesized row. If a row passes, ordinary positions contribute at most $1-(1-q^\star)^T$ to SER, and rare calibration-tail positions contribute at most $\delta T$.

\section{Sparse Asymmetric Coefficient Attack under OMP}
\label{app:omp-end-to-end-main}

This section composes the preceding results into the end-to-end guarantee used by the main attack chain. The constants are measurable properties of the selected anchors, the designed row, and the base/donor hidden-state distributions; no black-box SER calibration function is assumed.

\begin{assumption}[End-to-end sparse OMP attack setup]
\label{ass:end-to-end-omp-setup}
Fix a support $S\subseteq\mathcal{T}$ with $|S|=k$. The donor dictionary satisfies Assumption~\ref{ass:omp-coherence}, and the planted donor row satisfies Assumption~\ref{ass:omp-noisy-row}. Define
\[
\varepsilon_{\mathrm{rec}}
:=
\frac{\eta_\varepsilon}{\sigma_{\min}(\Phi_{d,S})}.
\]
Assume base alignment and donor innocuity:
\[
\|\Phi_{b,S}^\top\alpha-\mu_{\text{base}}\|_2\le \eta_b,
\qquad
\|U\Phi_{d,S}^\top\alpha\|_2\le \eta_d.
\]
Let
\[
M_d:=\|\Phi_{d,S}^\top\alpha\|_2+\eta_\varepsilon.
\]
\end{assumption}

\begin{assumption}[Base and donor hidden-state controls for SER]
\label{ass:end-to-end-hidden-state-controls}
For the attacked base, Assumption~\ref{ass:hidden-norm-control} holds with mean $\mu_h$, concentration radius $\xi_b$, failure probability $\zeta_b$, competitor row norm bound $R_b$, and vocabulary size $V_b$. Define
\[
\eta_h:=\|\mu_{\text{base}}-\mu_h\|_2.
\]
For the patched donor, let $h_t^{(d)}$ be the donor hidden state at a generation position under histories that have not yet emitted $\tau$. Let $P_{U^\perp}:=I-U^\top U$ after orthonormalizing the rows of $U$. Assume that, for every no-prior-$\tau$ history,
\[
\Pr\left[
\|Uh_t^{(d)}\|_2\le H_{d,U},\quad
\|P_{U^\perp}h_t^{(d)}\|_2\le H_{d,\perp},\quad
Z_d(h_t^{(d)})\ge Z_{d,0}
\mid \text{history}\right]
\ge
1-\zeta_d,
\]
where
\[
Z_d(h):=\sum_{v\in\mathcal{V}_d}\exp((w_v^{(d)})^\top h)
\]
is the natural donor logit partition excluding the breaker token. The constants $H_{d,U}$, $H_{d,\perp}$, $Z_{d,0}$, and $\zeta_d$ are measurable from the patched donor and the deployment prompt-position distribution.
\end{assumption}

\begin{theorem}[End-to-end sparse asymmetric coefficient attack under OMP]
\label{thm:end-to-end-omp-attack}
Under Assumptions~\ref{ass:end-to-end-omp-setup} and~\ref{ass:end-to-end-hidden-state-controls}, the victim's OMP-based transplant produces a base row $\widehat{x}_b=\Phi_b^\top\beta$ with the following properties.
\begin{enumerate}[leftmargin=*]
\item OMP recovers the support $S$ in exactly $k$ iterations, and
\[
\|\beta-\alpha\|_2\le \varepsilon_{\mathrm{rec}}.
\]
\item The transplanted base row satisfies
\[
\|\widehat{x}_b-\mu_{\text{base}}\|_2
\le
\eta_b+\|\Phi_{b,S}\|_{\mathrm{op}}\varepsilon_{\mathrm{rec}}.
\]
Consequently,
\[
\|\widehat{x}_b-\mu_h\|_2
\le
r_b
:=
\eta_h+\eta_b+\|\Phi_{b,S}\|_{\mathrm{op}}\varepsilon_{\mathrm{rec}}.
\]
\item The donor-side projection of the planted row satisfies
\[
\|Ux_d\|_2
\le
r_d
:=
\eta_d+\|U\|_{\mathrm{op}}\eta_\varepsilon.
\]
\item If $m_b(r_b)>0$, with $m_b$ from Lemma~\ref{lem:logit-margin-from-alignment}, define
\[
q_b
:=
\frac{1}{1+(V_b-1)\exp[-m_b(r_b)]}.
\]
Then the attacked-base length-$T$ sampling SER is lower-bounded by
\[
\operatorname{SER}_{\mathrm{base}}(\tau)
\ge
1-\left(1-(1-\zeta_b)q_b\right)^T.
\]
\item Define the donor breaker-logit upper bound
\[
L_d
:=
r_d H_{d,U}+M_d H_{d,\perp}
\]
and
\[
q_d
:=
\frac{\exp(L_d)}{Z_{d,0}+\exp(L_d)}.
\]
Then the patched-donor length-$T$ sampling SER is upper-bounded by
\[
\operatorname{SER}_{\mathrm{donor}}(\tau)
\le
1-(1-q_d)^T+\zeta_d T.
\]
The explicit asymmetry gap certified by the theorem is therefore at least
\[
\left[1-\left(1-(1-\zeta_b)q_b\right)^T\right]
-
\left[1-(1-q_d)^T+\zeta_d T\right].
\]
\end{enumerate}
\end{theorem}

\begin{proof}
The first claim is Theorem~\ref{thm:omp-coefficient-recovery}. The second claim is Corollary~\ref{cor:omp-base-reconstruction-error} with the design residual bounded by $\eta_b$ and $\|\Phi_{b,S}\|_{\mathrm{op}}$ used in place of $\|\Phi_b\|_{\mathrm{op}}$ on the recovered support. Adding $\|\mu_{\text{base}}-\mu_h\|_2=\eta_h$ gives the bound $r_b$.

For the donor projection,
\[
Ux_d=U\Phi_{d,S}^\top\alpha+U\varepsilon,
\]
so
\[
\|Ux_d\|_2
\le
\eta_d+\|U\|_{\mathrm{op}}\eta_\varepsilon
=r_d.
\]

For the base SER, Lemma~\ref{lem:logit-margin-from-alignment} applied with $w_\tau=\widehat{x}_b$ gives softmax probability at least $q_b$ on the concentration event, which has conditional probability at least $1-\zeta_b$ at each no-prior-$\tau$ step. Thus, conditional on not having emitted $\tau$ before step $t$, the probability of emitting $\tau$ at step $t$ is at least $(1-\zeta_b)q_b$. Iterating the probability of non-emission gives
\[
\Pr[\tau\text{ never appears in }T\text{ steps}]
\le
\left(1-(1-\zeta_b)q_b\right)^T,
\]
which proves the lower bound.

For the donor side, on the donor concentration event,
\[
x_d^\top h_t^{(d)}
=
(Ux_d)^\top Uh_t^{(d)}+(P_{U^\perp}x_d)^\top P_{U^\perp}h_t^{(d)}
\le
r_dH_{d,U}+\|x_d\|_2H_{d,\perp}.
\]
Since $\|x_d\|_2\le M_d$, this logit is at most $L_d$. The natural donor partition is at least $Z_{d,0}$ on the same event, so the breaker softmax probability is at most
\[
\frac{\exp(L_d)}{Z_{d,0}+\exp(L_d)}=q_d.
\]
Conditional on no tail event across the $T$ positions, the probability of at least one donor emission is at most $1-(1-q_d)^T$. The union bound gives probability at most $\zeta_d T$ for a tail event, proving the donor SER upper bound.
\end{proof}

\noindent\textbf{Intuition.}
The end-to-end chain has no hidden calibration oracle. Donor coherence and beta-min recover $\beta\approx\alpha$; conditioning of $\Phi_{d,S}$ controls the coefficient error; $\Phi_{b,S}$ passes this error to the base row; base hidden-state concentration converts row alignment into a softmax margin; and donor hidden-state concentration plus the partition lower bound converts donor innocuity into a per-step emission cap.
